\chapter{Music and Instrument Signals}

We have spent many chapters learning to take signals apart -- into sinusoids, into spectra, into
features. This chapter turns that machinery on one of the richest signal families humans have ever
made: the sounds of musical instruments, and the control language we use to drive them. There are
two reasons to care. The first is that music exercises nearly everything we have
built. A plucked string is a living demonstration of harmonics, decay, and timbre; a struck drum is
a lesson in inharmonicity; a chord is polyphony you can hear. The second reason is the goal of
this book: at the end of it we want to build a \emph{gesture controller} -- something worn on the
wrist and forearm that turns motion of the hand into music. To build that, we first have to know
exactly what ``music'' means as a signal and as a stream of control messages, and we have to reckon
with the tight timing budget a playable instrument imposes. By the end of this chapter you
should be able to look at a note and see four things at once: its pitch as a frequency, its timbre as
a spectrum, its life as an amplitude envelope, and its representation as a handful of bytes called
MIDI. That last picture is the target our controller will aim at.

\section{A note as a spectrum with a shape}
\label{ch29:note-is-spectrum}

Ask someone what the note A above middle C ``is'' and they will say $440$ Hz. That is true but tells
you little. A pure $440$ Hz sinusoid sounds like a hearing test, not like a violin, an oboe, or a
piano playing A. Every real instrument playing that note produces energy at $440$ Hz \emph{and} at a
whole ladder of higher frequencies, and it does so with an amplitude that rises and falls over the
life of the note. The pitch is one number; the actual signal is a spectrum with a time-varying shape.
The job of this section is to pull those two ideas -- the ladder (harmonics) and the shape over time
(the envelope) -- apart so we can study each on its own, and then put them back together.

\begin{intuitionbox}
A single instrument note is, to a first approximation, a sum of sinusoids whose frequencies are
\emph{integer multiples} of one fundamental, each with its own amplitude, the whole thing multiplied
by a slow envelope that swells and fades. Pitch lives in the spacing of the ladder. Timbre -- why a
trumpet and a clarinet at the same pitch sound different -- lives in the \emph{relative heights} of
the rungs. Attack and decay -- why a piano sounds nothing like an organ -- live in the envelope. Once
these three are separate in your mind, the rest of the chapter follows.
\end{intuitionbox}

Why integer multiples? The integer ratios follow from the physics
of the vibrating object. A string clamped at both ends, or an air column in a pipe, can only
support standing waves that ``fit'' the boundary -- a half wavelength, a full wavelength, three
halves, and so on. Those allowed modes have frequencies $f_0, 2f_0, 3f_0, \dots$, an arithmetic
ladder built on the lowest mode $f_0$. We met exactly this structure in the Fourier series: any
periodic signal of period $T_0 = 1/f_0$ is a sum of sinusoids at $f_0$ and its integer multiples. A
sustained instrument note is (nearly) periodic, so its spectrum is (nearly) a comb of spikes at the
harmonics.

\begin{connectionbox}
Mathematically, a sustained musical note is a direct realization of the Fourier series (Chapter~6). In the harmonic model of Section~\ref{ch29:note-is-spectrum},
a note is periodic with fundamental $f_0$; the Fourier-series coefficients $c_k$ are exactly the
complex amplitudes of the $k$-th harmonic. ``Timbre'' is just the magnitude profile $\abs{c_k}$
versus $k$. Everything we proved about Fourier series -- orthogonality of the harmonics, Parseval's
energy identity -- applies verbatim to the steady part of a musical tone.
\end{connectionbox}

\subsection*{Naming the ladder}

Some vocabulary, because the words trip people up. The \emph{fundamental} is $f_0$, the lowest mode,
and it is the one your ear reports as ``the pitch.'' The \emph{harmonics} are the whole ladder $f_0,
2f_0, 3f_0, \dots$; the $k$-th harmonic sits at $k f_0$. \emph{Overtones} are all the modes
\emph{above} the fundamental, so the first overtone is the second harmonic. (The off-by-one is an
unfortunate historical convention. We will say ``$k$-th harmonic'' and avoid ``overtone'' to stay safe.)

\begin{definition}[Harmonic series of a tone]
\label{ch29:def-harmonics}
A periodic tone of fundamental frequency $f_0$ has spectral energy concentrated at the frequencies
$f_k = k f_0$ for $k = 1, 2, 3, \dots$. The complex amplitude $c_k$ of the $k$-th harmonic sets both
its loudness ($\abs{c_k}$) and its phase ($\arg c_k$). The collection $\set{\abs{c_k}}$, viewed as a
shape over $k$, is the tone's \emph{harmonic spectrum} or, loosely, its \emph{timbre}.
\end{definition}

\begin{figure}[H]
\centering
\begin{tikzpicture}
\begin{axis}[
    width=12.8cm, height=5.4cm,
    xlabel={frequency (multiples of $f_0 = 220$ Hz)}, ylabel={amplitude},
    xmin=0, xmax=2400, ymin=0, ymax=1.1,
    axis lines=left, ytick={0,0.5,1},
    xtick={220,440,880,1320,1760,2200},
    xticklabels={$f_0$,$2f_0$,$4f_0$,$6f_0$,$8f_0$,$10f_0$},
    xticklabel style={font=\footnotesize},
    grid=both, grid style={gray!16},
]
\addplot[ycomb,sigblue,very thick,mark=*,mark options={scale=0.7}] coordinates {
  (220,1.0) (440,0.62) (660,0.45) (880,0.30) (1100,0.26)
  (1320,0.17) (1540,0.15) (1760,0.10) (1980,0.09) (2200,0.06)};
\end{axis}
\end{tikzpicture}
\caption{The harmonic comb of a sustained tone at $f_0 = 220$ Hz (the note A3). Energy appears only
at integer multiples of the fundamental; the falling heights are this instrument's timbre. Change the
spacing and you change the pitch; change the heights and you change which instrument it sounds like.}
\label{ch29:fig-comb}
\end{figure}

\subsection*{Timbre is the shape of the comb}

Two notes at the same pitch have the same comb \emph{spacing}; what differs is the \emph{profile} of
heights. A few archetypes worth carrying in your head:

\begin{itemize}
  \item A \textbf{flute} or a gently bowed string near the bridge can be almost a pure fundamental
    with a couple of weak harmonics -- a ``round,'' soft sound.
  \item A \textbf{clarinet} famously suppresses the even harmonics (its cylindrical bore closed at one
    end behaves like the odd-harmonic series), so the $3f_0, 5f_0, 7f_0$ rungs dominate. That hollow,
    woody quality is the even harmonics being missing.
  \item A \textbf{trumpet} or a bowed string near the bridge is rich in high harmonics -- many tall
    rungs far up the ladder -- which we hear as ``bright'' or ``buzzy.''
  \item A \textbf{sawtooth} (the staple of analog synths) contains \emph{every} harmonic with
    amplitude $1/k$; a \textbf{square} wave contains only odd harmonics, also $1/k$. These are not
    instruments but they are the reference shapes synthesis is built from.
\end{itemize}

\begin{figure}[H]
\centering
\begin{tikzpicture}
\begin{axis}[
    width=12.8cm, height=5.0cm,
    xlabel={time $t$ (one period)}, ylabel={pressure / level},
    xmin=0, xmax=1, ymin=-1.5, ymax=1.5,
    axis lines=left, samples=200, domain=0:1,
    legend style={at={(0.97,0.95)}, anchor=north east, font=\footnotesize, draw=gray!30, fill=white, fill opacity=0.9, text opacity=1, rounded corners=2pt}, legend cell align=left,
    every axis plot/.append style={very thick},
    grid=both, grid style={gray!16},
]
\addplot[sigblue] {sin(deg(2*pi*x))};
\addlegendentry{pure fundamental}
\addplot[specred,dashed] {sin(deg(2*pi*x)) + 0.5*sin(deg(4*pi*x)) + 0.333*sin(deg(6*pi*x)) + 0.25*sin(deg(8*pi*x)) + 0.2*sin(deg(10*pi*x)) + 0.166*sin(deg(12*pi*x))};
\addlegendentry{rich (sawtooth-like)}
\addplot[quotientorange,densely dotted] {sin(deg(2*pi*x)) + 0.333*sin(deg(6*pi*x)) + 0.2*sin(deg(10*pi*x))};
\addlegendentry{odd harmonics (hollow)}
\end{axis}
\end{tikzpicture}
\caption{Same pitch (same period), three timbres. The waveform shape in time is just a different way
of reading the harmonic profile: adding more, taller high harmonics sharpens the corners of the wave,
which the ear hears as a brighter tone.}
\label{ch29:fig-timbres}
\end{figure}

\begin{whybox}
Why does ``more high harmonics'' sound ``brighter''? Sharp features in time -- corners, jumps, fast
wiggles -- require high-frequency content to build, exactly as a square wave needs its endless odd
harmonics to make its vertical edges. The ear's cochlea is itself a bank of frequency-tuned filters
(Chapter~26), so it literally measures how much energy sits high up the ladder. A spectrum weighted
toward the top excites the high-frequency channels more, and the brain labels that ``bright.'' Timbre
perception is largely the ear reading the shape of the comb.
\end{whybox}

\section{The envelope: attack, decay, and why a piano is not an organ}
\label{ch29:envelope}

The sustained spectrum does not describe the whole sound, and for percussive instruments it matters less than the envelope.
Pluck a string and the sound is loudest at the very start and then fades; bow a string and you can
hold it forever; strike a piano key and you get a sharp onset followed by a long, slow decay that
never quite holds steady. That overall loudness-versus-time curve is the \emph{amplitude envelope},
and it carries a large share of the instrument's identity. The classic demonstration: record a piano note, cut
off the first tenth of a second -- the \emph{attack transient} -- and play the rest. It no longer
sounds like a piano at all; people often guess organ or flute. The steady harmonics are barely
changed; removing only the beginning is enough to lose the instrument's identity.

\begin{historybox}
This principle was established systematically by French composer-physicist Jean-Claude Risset
and computer-music pioneer Max Mathews at Bell Laboratories between 1964 and 1969. In the 19th
century, Hermann von Helmholtz's seminal acoustic treatise established the dogma that instrument
timbre was governed solely by its static Fourier harmonic recipe $\abs{c_k}$. But when early digital
synthesizers rendered notes with fixed harmonic proportions, the results sounded synthetic and
lifeless. Using Mathews's MUSIC IV software, Risset performed pitch-synchronous digital analysis
on real trumpet, piano, and string recordings, discovering that convincing acoustic identity
resides in the asynchronous micro-evolution of the attack transient: high-frequency harmonics
often rise late and decay earliest. This work moved beyond the static Fourier view of musical
tone, laid the groundwork for dynamic additive and physical modeling synthesis, and showed that
timbre is inherently time-varying. An active open frontier is real-time neural modeling of non-linear
chaotic acoustic transients (such as the bite of a cello bow or the multiphonics of woodwinds)
under microsecond performance latencies.
\end{historybox}

\subsection*{The ADSR model}

The standard parametric description of an envelope, borrowed from synthesizers, breaks the note into four
phases, the \emph{ADSR} model:

\begin{itemize}
  \item \textbf{Attack} -- the rapid rise from silence to peak level when the note begins. For a
    plucked or struck note this is milliseconds; for a slowly bowed note it can be hundreds of
    milliseconds. This is the \emph{attack transient}, and it is where most of the instrument's
    fingerprint lives.
  \item \textbf{Decay} -- the drop from the peak down to a steady level after the initial onset.
  \item \textbf{Sustain} -- the held level while the note is kept active (a key held down, a bow kept
    moving). Note this is a \emph{level}, not a time; a piano has essentially no true sustain, its
    sound decaying throughout.
  \item \textbf{Release} -- the fade to silence after the note is let go (key up).
\end{itemize}

\begin{figure}[H]
\centering
\begin{tikzpicture}[scale=1.0]
\draw[->] (0,0) -- (10.5,0) node[right]{time};
\draw[->] (0,0) -- (0,3.5) node[left, font=\small]{amplitude};
\draw[very thick,sigblue]
  (0,0) -- (1.4,3.0)
  -- (3.0,1.8)
  -- (6.6,1.6)
  -- (8.6,0);
\foreach \x in {1.4,3.0,6.6,8.6} {
  \draw[gray!55,dashed] (\x,0) -- (\x,3.1);
}
\node at (0.7,-0.35){\small attack};
\node at (2.2,-0.35){\small decay};
\node at (4.8,-0.35){\small sustain};
\node at (7.6,-0.35){\small release};
\draw[<->,gray!70] (0.05,3.25) -- (1.35,3.25);
\node[gray!70, font=\scriptsize] at (0.7,3.5){transient};
\node[sigblue] at (6.7,1.95){\scriptsize key released};
\end{tikzpicture}
\caption{The ADSR amplitude envelope. The note's identity is concentrated in the steep attack
transient at the left; the sustain is a held level, not a duration, and the release begins when the
player lets go. A piano has a tall attack and essentially no flat sustain (the whole tail is decay);
an organ is nearly all sustain with negligible attack and release.}
\label{ch29:fig-adsr}
\end{figure}

\begin{happeningbox}
What is really happening when you ``hear the attack''? Two things at once. First, the broadband
\emph{transient} -- the click of the hammer, the scrape of the bow catching, the chiff of air -- is a
short, sharp event, and short sharp events are spectrally wide (the time--bandwidth trade-off again).
Second, during the attack the harmonics are not yet in their steady balance: they ramp up at
different rates and the pitch itself may be momentarily unstable. The ear is highly sensitive to
these first tens of milliseconds because in the natural world the onset of a sound is where the
information about ``what made it'' is densest. This is also exactly why a controller that lags the
onset feels dead, a point we return to in Section~\ref{ch29:latency}.
\end{happeningbox}

\section{Inharmonicity: when the ladder bends}
\label{ch29:inharmonicity}

The integer-multiple story is an idealization. Real strings have stiffness, so their higher modes are
pulled slightly \emph{sharp} of the perfect ladder: the $k$-th partial of a real piano string sits
not at $k f_0$ but at roughly
\[
  f_k \;=\; k f_0 \sqrt{1 + B\,k^2},
\]
where $B$ is a small \emph{inharmonicity coefficient} that grows with string stiffness and is larger
for the short, thick strings of the bass. This is why a piano is tuned with ``stretched'' octaves:
the top notes are tuned a touch sharp and the bottom notes a touch flat so that the actual partials
of one note line up with the actual partials of the next. A perfectly mathematically tuned piano
sounds out of tune.

\begin{trapbox}
Do not assume a pitch detector that hunts for perfectly integer-spaced peaks will work on real
instruments. Piano, marimba, and especially bells violate the integer ladder badly. A bell's partials
can be wildly inharmonic; the ``pitch'' you hear is sometimes a \emph{missing fundamental} that the
brain infers from the spacing of partials that are not even harmonic. Pitch is an auditory
percept of periodic organization, rather than simply the tallest spike in a Fourier magnitude plot.
\end{trapbox}

\begin{applicationbox}
Drums and cymbals abandon the ladder entirely. A circular drumhead's modes are set by the zeros of
Bessel functions, not by integers, so the partials are inharmonic and you perceive a sense of pitch
only weakly (more for tom and timpani, almost none for a snare). A cymbal is closer to filtered noise
than to a tone. This matters for the controller: \emph{percussive} gestures are far easier to map
than \emph{pitched} ones, because a drum hit only needs an onset and a velocity, not a stable
frequency. The first instrument our wrist controller will convincingly play is a drum kit.
\end{applicationbox}

\section{Pitch and tuning: from frequency to note number}
\label{ch29:pitch-tuning}

To turn ``a frequency'' into ``a note'' we need the map between hertz and the names musicians use. The
modern Western system is \emph{twelve-tone equal temperament}: the octave (a doubling of frequency)
is divided into twelve equal steps called \emph{semitones}, each a frequency ratio of $2^{1/12}
\approx 1.0595$. Everything is anchored to a reference, conventionally A4 $= 440$ Hz.

\begin{definition}[Equal-tempered frequency]
\label{ch29:def-et}
With A4 $= 440$ Hz as the anchor, the frequency $n$ semitones away is
\[
  f \;=\; 440 \cdot 2^{\,n/12} \quad\text{Hz}.
\]
Conversely, the number of semitones above A4 of a frequency $f$ is $n = 12\,\log_2(f/440)$, and
fractional $n$ measures how far the note is mistuned. One hundredth of a semitone is a \emph{cent}, so
the mistuning in cents is $1200\,\log_2(f/f_{\text{target}})$.
\end{definition}

\begin{example}[How far apart are adjacent notes?]
\label{ch29:ex-spacing}
One semitone up from $440$ Hz is $440\cdot 2^{1/12} \approx 466.2$ Hz, a gap of about $26$ Hz. But
one semitone up from $110$ Hz (A2) is $110\cdot 2^{1/12}\approx 116.5$ Hz, a gap of only $6.5$ Hz. The
\emph{ratio} is constant ($2^{1/12}$); the \emph{absolute} gap shrinks fourfold every octave you go
down. A pitch detector therefore needs much finer frequency resolution in the bass than in the
treble -- one of the reasons short DFT windows struggle with low notes (Chapter~10), and a reason the
constant-Q transform, which spaces its bins logarithmically, is a natural fit for music.
\end{example}

\begin{connectionbox}
The logarithm in Definition~\ref{ch29:def-et} is why music is a logarithmic phenomenon and why the
right frequency axis for music is $\log f$. On a log axis the twelve semitones are \emph{evenly}
spaced and every octave is the same width, as in Figure~\ref{ch29:fig-logaxis}. This is precisely the constant-Q / wavelet philosophy of
Chapter~24: constant \emph{ratio} bandwidth rather than constant \emph{absolute} bandwidth. The ear
does the same thing -- pitch perception is roughly logarithmic in frequency.
\end{connectionbox}

\begin{figure}[H]
\centering
\begin{tikzpicture}
\begin{axis}[
    width=12.8cm, height=4.8cm,
    xlabel={frequency (Hz)}, ylabel={}, ytick=\empty,
    xmode=log, log basis x=2,
    xmin=100, xmax=1900, ymin=0, ymax=1.15,
    axis lines=left,
    xtick={110,220,440,880,1760},
    xticklabels={110 (A2),220 (A3),440 (A4),880 (A5),1760 (A6)},
    xticklabel style={font=\scriptsize},
]
% the twelve semitones across the A2--A3 octave as a comb of equal log-spacing
\addplot[ycomb,sigblue,thick,mark=*,mark options={scale=0.5}] coordinates {
  (110,1) (116.54,1) (123.47,1) (130.81,1) (138.59,1) (146.83,1)
  (155.56,1) (164.81,1) (174.61,1) (185.00,1) (196.00,1) (207.65,1) (220,1)};
% the octave A's as taller red marks
\addplot[ycomb,specred,very thick] coordinates {(110,1.1)(220,1.1)(440,1.1)(880,1.1)(1760,1.1)};
\end{axis}
\end{tikzpicture}
\caption{The equal-tempered octave A2--A3 on a $\log_2$ frequency axis. The blue comb marks the twelve
semitones from $110$ to $220$ Hz; on this axis they are \emph{evenly} spaced. The taller red marks are
successive A's ($110, 220, 440, 880, 1760$ Hz), and every octave A2--A3--A4--A5 spans the same
horizontal width even though the absolute hertz gaps double each octave. This even spacing is why a log
axis is the natural one for music.}
\label{ch29:fig-logaxis}
\end{figure}

Finding $f_0$ from a recorded note is its own small art, because of the missing-fundamental problem
and the inharmonicity above. Two families of method dominate.

\textbf{Spectral.} Take a DFT, find the harmonic comb, and infer $f_0$ from the spacing rather than
from any single peak (a \emph{harmonic product spectrum} multiplies downsampled copies of the
magnitude spectrum so the true $f_0$ reinforces across harmonics). Robust to a weak or missing
fundamental; limited in the bass by the DFT's resolution.

\textbf{Time-domain autocorrelation.} A periodic signal of period $T_0$ correlates strongly with
itself shifted by $T_0$. So compute the autocorrelation and find the first strong peak away from zero
lag; its lag is $T_0$ and $f_0 = 1/T_0$. The widely used \emph{YIN} algorithm is a refined version
using a difference function instead of a raw correlation to avoid octave errors. This works even when
the fundamental is weak, because periodicity in time survives a missing fundamental.

\begin{lstlisting}[language=Python,caption={Autocorrelation pitch estimate (the idea, not production code).}]
import numpy as np

def estimate_f0(x, fs, fmin=50, fmax=1000):
    x = x - np.mean(x)
    corr = np.correlate(x, x, mode='full')[len(x)-1:]   # nonneg lags
    lo, hi = int(fs/fmax), int(fs/fmin)                 # lag search window
    lag = lo + int(np.argmax(corr[lo:hi]))             # first strong peak
    return fs / lag                                     # f0 in Hz
\end{lstlisting}

\begin{walkbox}
Suppose you sample at $f_s = 8000$ Hz and the autocorrelation peak lands at lag $182$ samples. What
note is it? The period is $182/8000 = 22.75$ ms, so $f_0 = 8000/182 \approx 43.96$ Hz. Then $n =
12\log_2(43.96/440) \approx -39.9$ semitones below A4, i.e. essentially F1. Notice the resolution
problem: the next available integer lag, $181$, gives $44.2$ Hz -- already $10$ cents away. In the
bass, single-sample lag steps are coarse. Real pitch trackers interpolate between lags (parabolic
fit) to do better.
\end{walkbox}

\section{Polyphony: more than one note at once}
\label{ch29:polyphony}

So far, one note. Real music is mostly \emph{polyphonic}: chords on a piano, several strings in an
orchestra, a singer over a guitar. In the time domain this is simply addition -- the air pressure is
the sum of every source -- so the spectrum is the sum of every comb. That sounds easy and it is, going
\emph{forward}: to synthesize a chord you add the notes. Going \emph{backward} -- given a mixed
recording, recover which notes are playing -- is one of the genuinely hard problems in audio, called
\emph{multi-pitch estimation} or automatic music transcription.

\begin{figure}[H]
\centering
\begin{tikzpicture}
\begin{axis}[
    width=12.8cm, height=5.2cm,
    xlabel={frequency (Hz)}, ylabel={amplitude},
    xmin=0, xmax=1500, ymin=0, ymax=1.1,
    axis lines=left, ytick={0,0.5,1},
    grid=both, grid style={gray!16},
]
% C4 ~262 Hz comb (blue)
\addplot[ycomb,sigblue,very thick] coordinates {
  (262,1.0) (524,0.55) (786,0.40) (1048,0.25) (1310,0.18)};
% E4 ~330 Hz comb (red)
\addplot[ycomb,specred,very thick] coordinates {
  (330,0.9) (660,0.5) (990,0.34) (1320,0.22)};
% G4 ~392 Hz comb (orange)
\addplot[ycomb,quotientorange,very thick] coordinates {
  (392,0.85) (784,0.47) (1176,0.30)};
\end{axis}
\end{tikzpicture}
\caption{Three overlapping harmonic combs -- a C major chord (C4, E4, G4). The combined spectrum is
just the union of the three; a transcription algorithm must work backward to separate them. Note the
near-collision around $784$--$786$ Hz: the third harmonic of C4 and the second harmonic of G4 almost
coincide (C and G are a fifth apart), which is exactly \emph{why} that interval sounds consonant -- and
exactly why the two notes are hard to tell apart automatically.}
\label{ch29:fig-chord}
\end{figure}

\begin{whybox}
Why is multi-pitch hard when summing is so easy? Because harmonics \emph{collide}. Consonant
intervals are consonant precisely because their harmonics overlap: an octave shares every other
harmonic, a fifth shares one in three. When two notes share spectral energy, no amount of looking at
the magnitude spectrum can cleanly say how much belongs to each. The problem is ill-posed in the same
way as ``two numbers sum to $10$, what are they?'' Solving it requires priors -- about how many notes
are likely, about instrument timbres, about musical structure -- which is why modern transcription
leans on machine learning (Chapter~39). For our controller we sidestep this entirely: we
\emph{generate} polyphony from gestures rather than trying to \emph{hear} it.
\end{whybox}

\section{MIDI: music as a stream of bytes}
\label{ch29:midi}

Everything above describes music as a \emph{signal} -- pressure versus time. But there is a second,
quite different representation that our controller will actually produce: \emph{MIDI}, the Musical
Instrument Digital Interface, a standard from 1983 for sending performance \emph{events} between
devices. MIDI does not carry audio. It carries instructions: ``press this key this hard now,''
``release it,'' ``bend the pitch,'' ``move this knob.'' A separate \emph{synthesizer} or
\emph{sampler} turns those instructions into sound. This split -- events here, sound there -- is the
most important idea for the capstone, because a wrist sensor can plausibly produce
\emph{events} even when it could never produce convincing \emph{audio}.

\begin{applicationbox}
MIDI was established in 1982--83 by Dave Smith (Sequential Circuits) and Ikutaro Kakehashi (Roland)
deliberately as an open, royalty-free standard so synthesizers from competing manufacturers could
interoperate. Designed for the slow microcontrollers and inexpensive optocoupled serial hardware
of that era ($31{,}250$ baud, unidirectional, asynchronous bytes), its frugality allowed the
standard to endure four decades virtually unchanged. While MIDI 2.0 (2020) introduces
32-bit resolution and bidirectional capability,
classic MIDI 1.0 remains the lightweight protocol of choice for embedded wearable instruments.
\end{applicationbox}

\subsection*{Anatomy of a message}

A classic MIDI message is one \emph{status} byte followed by one or two \emph{data} bytes. The status
byte has its high bit set ($\ge 128$) and encodes both the message type and a 4-bit channel ($0$--$15$,
so up to 16 instruments on one wire). Data bytes have their high bit clear, so each carries 7 bits, a
value $0$--$127$. The handful you must know:

\begin{itemize}
  \item \textbf{Note On} -- status \texttt{0x90 | channel}, then \emph{note number} ($0$--$127$) and
    \emph{velocity} ($0$--$127$). Start a note. Velocity is how hard/fast the key was struck -- usually
    mapped to loudness and often to brightness too.
  \item \textbf{Note Off} -- status \texttt{0x80 | channel}, then note number and release velocity.
    Stop a note. (A Note On with velocity $0$ is, by convention, also a Note Off -- a trick that lets
    a stream of running Note Ons turn notes off cheaply.)
  \item \textbf{Control Change (CC)} -- status \texttt{0xB0 | channel}, then a \emph{controller
    number} ($0$--$127$) and a \emph{value} ($0$--$127$). The general-purpose continuous knob: CC~1 is
    the modulation wheel, CC~7 is channel volume, CC~11 is expression, CC~64 is the sustain pedal.
  \item \textbf{Pitch Bend} -- status \texttt{0xE0 | channel}, then two data bytes forming a 14-bit
    value ($0$--$16383$, center $8192$). A high-resolution continuous slide in pitch.
  \item \textbf{Program Change} -- status \texttt{0xC0 | channel}, then one byte: choose which
    instrument ``patch'' this channel plays.
\end{itemize}

\begin{definition}[MIDI note number]
\label{ch29:def-midinote}
MIDI numbers the equal-tempered keyboard so that A4 $= 440$ Hz is note $69$, and each semitone is one
unit. The frequency of MIDI note $m$ is
\[
  f(m) \;=\; 440 \cdot 2^{(m-69)/12} \quad\text{Hz},
\]
and the nearest MIDI note to a measured frequency $f$ is $m = \operatorname{round}\!\big(69 +
12\log_2(f/440)\big)$. Middle C (C4) is note $60$; the standard 88-key piano spans notes $21$ (A0) to
$108$ (C8).
\end{definition}

\begin{example}[Bytes for a note]
\label{ch29:ex-bytes}
To play middle C at medium-strong velocity on channel 1, then release it, you send six bytes:
\[
  \underbrace{\texttt{0x90}}_{\text{Note On, ch 1}}\;
  \underbrace{\texttt{0x3C}}_{60\,=\,\text{C4}}\;
  \underbrace{\texttt{0x64}}_{100\,=\,\text{velocity}}
  \qquad\text{then later}\qquad
  \underbrace{\texttt{0x80}}_{\text{Note Off}}\;
  \underbrace{\texttt{0x3C}}_{60}\;
  \underbrace{\texttt{0x00}}_{\text{rel. vel.}}.
\]
Three bytes to start a note, three to stop it. At the MIDI baud rate each byte takes about $0.32$ ms,
so a Note On is on the wire in roughly $1$ ms -- a number we will need in
Section~\ref{ch29:latency}.
\end{example}

\begin{trapbox}
MIDI is \emph{not} sound and a MIDI note number is \emph{not} a frequency in the analog sense. Note
$60$ is exactly C4 -- you cannot send ``a little sharp'' with a note number alone; for continuous
pitch you must add Pitch Bend (14-bit) on top. Likewise, everything in classic MIDI is quantized to
$7$ bits ($128$ levels). Velocity has only $128$ steps; CC has only $128$ steps. For a gesture
controller streaming a smoothly moving hand, that coarse quantization can produce audible ``zipper''
steps -- a real design constraint we will smooth around, or escape with MIDI 2.0's higher resolution.
\end{trapbox}

\begin{applicationbox}
Because MIDI is just events, the \emph{same} performance can drive a grand piano, a string section, or
a drum machine -- you only change the synth on the other end, or the Program Change you send. Channel
10 is reserved by the General MIDI convention for percussion, where the note \emph{number} chooses the
drum (note $36$ = kick, $38$ = snare, $42$ = closed hi-hat) instead of a pitch. This is why drums are
our easiest first target: a gesture becomes a single Note On to channel 10 with the right number, and
velocity becomes how hard you hit.
\end{applicationbox}

\subsection*{Continuous control versus discrete events}

Notice the two flavors of message, because they map onto two flavors of gesture. \emph{Note On/Off}
are \emph{discrete, timed events} -- a tap, a strike, a trigger. \emph{CC and Pitch Bend} are
\emph{continuous streams} -- a tilt, a sweep, a squeeze, sampled and sent maybe $50$--$200$ times a
second. A good controller uses both: a sharp wrist flick fires a Note On (a drum hit, a plucked
note); a slow forearm rotation rides a CC (a filter sweep, a swell in volume). Mapping which motions
become events and which become streams is the core design problem of the capstone.

\begin{lstlisting}[language=Python,caption={Turn a measured pitch and a gesture intensity into MIDI bytes.}]
def freq_to_note(f):
    return round(69 + 12 * math.log2(f / 440.0))

def note_on(note, velocity, channel=0):
    return bytes([0x90 | channel, note & 0x7F, velocity & 0x7F])

def note_off(note, channel=0):
    return bytes([0x80 | channel, note & 0x7F, 0])

def control_change(cc, value, channel=0):
    return bytes([0xB0 | channel, cc & 0x7F, value & 0x7F])

# Example: a flick of intensity 0..1 fires a snare hit on the drum channel (9 == "channel 10")
def flick_to_snare(intensity):
    vel = max(1, min(127, int(intensity * 127)))
    return note_on(38, vel, channel=9)
\end{lstlisting}

\section{Latency: the unforgiving budget of a playable instrument}
\label{ch29:latency}

One constraint shapes every design choice in the capstone. An instrument must
respond to the player \emph{fast enough that the response feels simultaneous with the action}, so
that the delay is imperceptible. The human perceptual threshold for action-to-sound delay is
tight, and missing it makes an instrument feel broken rather than merely
slow, because your hand and your ear disagree about when the note happened.

\begin{focusbox}
The numbers, roughly, from the perception and music-technology literature. Around \textbf{10 ms} of
action-to-sound latency is the comfortable target -- below it almost nobody notices. Around
\textbf{20--30 ms} skilled players begin to feel a sluggishness, especially on percussive
attacks. Past \textbf{40--50 ms} the delay is plainly audible to most people and rhythmic playing
falls apart. For comparison, standing $3.4$ m from your own amplifier already costs $10$ ms in sound
travel time, which is why drummers are so sensitive to it. The whole signal chain -- sense, process,
transmit, synthesize, output -- must fit inside this budget.
\end{focusbox}

Why is the attack so much less forgiving than a held note? Because, as Section~\ref{ch29:envelope}
argued, the onset is where the ear's clock is set. A held swell can drift in by a few extra
milliseconds and nobody minds; a drum hit that lands $30$ ms after the flick feels rubbery and wrong.
Percussive, transient-heavy gestures are simultaneously the \emph{easiest} to detect (a sharp event
is easy to spot) and the \emph{least} forgiving of delay. That tension is the central engineering
problem of the controller.

\begin{figure}[H]
\centering
\begin{tikzpicture}[font=\small, >=Stealth]
\tikzset{
  sens/.style={draw=blue!70!black, fill=sourceblue!10, rounded corners=2pt, minimum height=11mm, minimum width=22mm, align=center, thick},
  buf/.style={draw=quotientorange!80!black, fill=quotientorange!12, rounded corners=2pt, minimum height=11mm, minimum width=22mm, align=center, thick},
  proc/.style={draw=targetgreen!70!black, fill=targetgreen!10, rounded corners=2pt, minimum height=11mm, minimum width=22mm, align=center, thick},
  midi/.style={draw=structurepurple!70!black, fill=structurepurple!12, rounded corners=2pt, minimum height=11mm, minimum width=20mm, align=center, thick},
  synth/.style={draw=specred!70!black, fill=specred!10, rounded corners=2pt, minimum height=11mm, minimum width=22mm, align=center, thick}
}

\node[sens] (s) {Acquisition\\ IMU / EMG\\ {\footnotesize 1--3\,ms}};
\node[buf] (b) [right=6mm of s] {Buffer \&\\ Window\\ {\footnotesize 2--6\,ms}};
\node[proc] (d) [right=6mm of b] {Detect \&\\ Classify\\ {\footnotesize 1--4\,ms}};
\node[midi] (m) [right=6mm of d] {Emit MIDI\\ Event\\ {\footnotesize $\sim$1\,ms}};
\node[synth] (y) [right=6mm of m] {Audio Engine\\ Synthesis\\ {\footnotesize 3--6\,ms}};

\draw[->, thick] (s) -- (b);
\draw[->, thick] (b) -- (d);
\draw[->, thick] (d) -- (m);
\draw[->, thick] (m) -- (y);

% Overall budget bracket
\coordinate (bl) at ($(s.south west)+(0,-0.4)$);
\coordinate (br) at ($(y.south east)+(0,-0.4)$);
\draw[{Bar[width=5pt]}-{Bar[width=5pt]}, gray!80, thick] (bl) -- (br);
\node[below=5pt, font=\footnotesize\bfseries, gray!80] at ($(bl)!0.5!(br)$) {Total Latency Budget $\le$ 15\,ms (Perceptual Simultaneity Threshold)};
\end{tikzpicture}
\caption{The end-to-end action-to-sound latency budget for a playable gesture controller. Every stage between muscle initiation and speaker output chips away at the strict $\sim$15\,ms envelope. Buffering is typically the largest component: triggering immediately on rising edges rather than waiting for complete gesture windows keeps the system responsive and tactile.}
\label{ch29:fig-latency}
\end{figure}

\begin{trapbox}
The latency villain is almost always \emph{buffering}, not computation. To analyze a gesture you
collect a window of samples; a $256$-sample window at $1$ kHz is already $256$ ms -- catastrophic. You
cannot wait for a long window and also be fast. This is the same time--frequency uncertainty we have
met throughout the book: a longer window buys cleaner frequency/feature estimates but costs you
exactly that much delay. Onset detection that triggers on the \emph{leading edge} of a gesture,
rather than waiting to see the whole thing, is how playable controllers avoid this trap: they accept
occasional false triggers in exchange for low latency.
\end{trapbox}

\begin{figure}[H]
\centering
\begin{tikzpicture}
\begin{axis}[
    width=12.8cm, height=5.2cm,
    xlabel={time (ms)}, ylabel={jerk magnitude (a.u.)},
    xmin=0, xmax=120, ymin=0, ymax=1.25,
    axis lines=left, samples=160, domain=0:120,
    grid=both, grid style={gray!16},
    legend pos=north east, legend cell align=left,
]
% a flick: quiet, then a sharp burst of jerk that rises and falls
\addplot[sigblue,very thick,smooth] {0.06 + 1.05*exp(-((x-45)^2)/40)};
\addlegendentry{jerk}
\addplot[specred,dashed,thick] coordinates {(0,0.35) (120,0.35)};
\addlegendentry{threshold}
% rising-edge crossing: vertical mark where jerk first crosses the threshold
\addplot[gray!70,densely dotted,very thick] coordinates {(38.5,0) (38.5,0.35)};
\addlegendentry{fire Note On}
\end{axis}
\end{tikzpicture}
\caption{A wrist flick seen as jerk magnitude (derivative of acceleration) from the IMU. The Note On
fires the instant the rising edge crosses the threshold, near $t=38$ ms here -- well before the peak at
$t=45$ ms. Waiting for the peak, or for a full analysis window, would add the difference straight to
the latency budget. Triggering on the leading edge is what keeps a percussive gesture feeling
instant.}
\label{ch29:fig-onset}
\end{figure}

\begin{buildbox}
Concrete budget for a wrist controller streaming MIDI over BLE to a phone synth:
\begin{itemize}
  \item \textbf{Sense.} Sample the IMU at $\ge 200$ Hz (gyro/accel), EMG/FMG at $\ge 1$ kHz. One
    sample period is $1$--$5$ ms. Use the sensor's own FIFO/interrupt, do not poll lazily.
  \item \textbf{Detect.} Trigger Note On on the \emph{rising edge} of a feature (jerk magnitude, EMG
    envelope crossing a threshold), as in Figure~\ref{ch29:fig-onset}. Keep the decision window
    short -- tens of samples -- so detection adds only a few ms.
  \item \textbf{Transmit.} Classic serial MIDI puts a $3$-byte Note On on the wire in $\sim$1 ms.
    \textbf{Bluetooth LE MIDI is the trap}: the BLE connection interval (often $7.5$--$15$ ms, and on
    some phones much worse) can dominate your entire budget. Request the shortest connection interval
    the link allows, and prefer a wired/USB path when you need true tightness.
  \item \textbf{Synthesize.} A phone or laptop softsynth with a small audio buffer adds $3$--$10$ ms;
    on Android use the low-latency audio path (AAudio), on a Pi use a small ALSA period.
\end{itemize}
Measure it, do not trust the spec sheet: tap the sensor against a mic, record both the tap transient
and the synth output on one track, and read the gap. Anything under $\sim$15 ms feels immediate.
\end{buildbox}

\begin{connectionbox}
This section relies on results from earlier chapters. The window-length-versus-delay
trade-off is the time--frequency uncertainty principle of Chapters~7 and~10. ``Trigger on the rising
edge'' is onset/edge detection -- a high-pass or derivative filter (Chapter~16) followed by a
threshold. And ``the BLE interval dominates'' is just sampling theory applied to the transport: your
update rate is capped by how often the channel actually delivers, no matter how fast you compute.
\end{connectionbox}

\section{Setting the musical target for the capstone}
\label{ch29:target}

We can now state precisely what the gesture controller is trying to do, in the vocabulary of this
chapter. It is a device that converts forearm and wrist motion into a stream of MIDI events --
discrete Note On/Off for strikes and plucks, continuous CC and Pitch Bend for sweeps and bends -- fast
enough to feel simultaneous with the gesture, which a downstream synth turns into sound. That framing
already determines the hardware design and its practical limits.

\begin{goalbox}
\textbf{The target, and the limits that go with it.} Our controller's job is to emit musically
useful MIDI from wrist/forearm sensing within the latency budget of Section~\ref{ch29:latency}. What
that realistically allows, and what it does not:
\begin{itemize}
  \item \textbf{Strongly within reach: discrete, percussive triggers.} A wrist-worn IMU detects sharp
    motion -- taps, flicks, strikes -- very well. Map each to a Note On (a drum on channel~10, a
    plucked note), with gesture intensity to velocity. This is the first thing to build because onset
    detection is robust and drums tolerate the most timbral slop.
  \item \textbf{Within reach: continuous macro-control.} Forearm orientation from sensor fusion
    (Chapter~36) drives CC and Pitch Bend smoothly -- tilt for a filter sweep, rotation for volume, a
    held pose for a sustained swell. Watch the 7-bit ``zipper'' quantization of
    Section~\ref{ch29:midi}; smooth and rate-limit the stream.
  \item \textbf{Reachable only with the right sensor: coarse hand states.} A few \emph{discrete} hand
    configurations -- fist vs.\ open, a pinch -- can be read from \emph{EMG or force-myography} on the
    forearm, and mapped to mode switches or note selection. This requires electromyographic
    muscle sensing rather than inertial motion alone.
  \item \textbf{Out of reach, and we will say so plainly: continuous per-finger pitch.} A wrist-worn
    IMU \emph{cannot} recover continuous individual finger joint angles -- it senses the whole forearm,
    not the fingers. EMG/FMG can distinguish a handful of coarse gestures, not a precise continuous
    finger pose. Genuinely continuous finger tracking needs sensors \emph{on the fingers} (rings,
    gloves) or external cameras. So a wrist device is a fine \emph{drummer} and \emph{macro-controller}
    and a poor \emph{pianist}; we design for what the sensors can actually deliver, not for the demo we
    wish we could give.
\end{itemize}
\end{goalbox}

\begin{aestheticbox}
The reason to lay out the musical target now, chapters before we wire anything, is that the music
defines the spec. ``Playable'' has a concrete meaning: $\lesssim$ 15 ms of latency, $128$ velocity
levels, an integer note number plus a 14-bit bend, and an onset detector that fires on a leading edge.
Every later chapter -- IMUs, calibration, fusion, EMG/FMG, gesture recognition -- is in service of
hitting that spec. Knowing what a note \emph{is}, as both a signal and a message, is what lets us
judge whether the hardware we build is actually an instrument or merely a sensor that beeps.
\end{aestheticbox}

With musical instrument acoustics, envelopes, and latency budgets established, we
bring Part~VI to a close. In Part~VII, we move from one-dimensional temporal waveforms
to two-dimensional spatial signals: images. Chapter~30 explores how images are sampled,
represented, and transformed as 2D arrays.

\section{Exercises}
\label{ch29:exercises}

\begin{exercise}
A sustained oboe note has its strongest energy not in the fundamental but in the second and third
harmonics. (a) Sketch a plausible harmonic comb for fundamental $f_0 = 440$ Hz with this property.
(b) Explain, using Section~\ref{ch29:pitch-tuning}'s missing-fundamental idea, why you would still
hear the pitch as $440$ Hz even if you filtered the fundamental out entirely.
\end{exercise}

\begin{exercise}
Compute the frequencies of the first five harmonics of a piano note with $f_0 = 110$ Hz, first under
the ideal integer ladder and then under the inharmonic model $f_k = k f_0\sqrt{1 + B k^2}$ with $B =
0.0004$. By how many cents is the fifth partial sharp of $5 f_0$? (Recall cents $= 1200\log_2$ of the
ratio.)
\end{exercise}

\begin{exercise}
Using Definition~\ref{ch29:def-midinote}, find the MIDI note number nearest to (a) $261.6$ Hz, (b)
$1000$ Hz, (c) $55$ Hz. For each, report the residual mistuning in cents between the measured frequency
and the exact equal-tempered frequency of that note.
\end{exercise}

\begin{exercise}
Write the exact byte sequence (in hex) to: play A4 at velocity $80$ on channel~1, then while it is
held send a modulation-wheel CC of value $64$, then release the note. Use the message formats in
Section~\ref{ch29:midi}.
\end{exercise}

\begin{exercise}
A controller samples a gesture feature at $f_s = 500$ Hz and detects an onset by waiting for a full
$N$-sample window before deciding. (a) For $N = 8, 32, 128$, what latency does the buffering alone add?
(b) Which of these fit inside a $15$ ms total budget, leaving room for the other stages of
Figure~\ref{ch29:fig-latency}? (c) Explain why ``trigger on the rising edge'' beats ``wait for the
whole window.''
\end{exercise}

\begin{exercise}
A C major chord plays C4 ($\approx 262$ Hz), E4 ($\approx 330$ Hz), G4 ($\approx 392$ Hz). List the
first four harmonics of each and find every pair of harmonics from \emph{different} notes that fall
within $5$ Hz of one another. Relate your answer to the observation earlier in this chapter that acoustic consonance and overlapping harmonic spectra stem from the same small-integer frequency ratios.
\end{exercise}

\begin{exercise}
Classic MIDI velocity and CC are $7$-bit ($0$--$127$). Suppose a forearm-tilt gesture sweeps a CC
linearly over $2$ seconds and you send updates at $100$ Hz. (a) How many distinct CC values can the
sweep pass through? (b) If the synth maps CC to a filter cutoff over $5$ octaves, how many cents does
one CC step correspond to, and is that audible? (c) Suggest two ways to hide the steps.
\end{exercise}

\begin{exercise}
The Risset/Mathews experiment (Section~\ref{ch29:envelope}) showed that removing a note's attack
transient destroys its identity. Design a small experiment with free software (e.g.\ record a single
piano note, a single flute note) to test this: describe exactly which $\sim$50 ms you would delete,
what you would play to listeners, and what result would confirm that timbre lives in the attack.
\end{exercise}

\begin{exercise}
Explain, in two or three sentences each and citing the relevant claim in the chapter, why each of the
following capstone goals is realistic or not for a \emph{wrist-only IMU}: (a) triggering a snare drum
on a flick; (b) a continuous volume swell from forearm rotation; (c) recovering the continuous bend
angle of the index finger; (d) distinguishing a closed fist from an open hand.
\end{exercise}

\begin{exercise}
(Open-ended build.) Sketch the block diagram of a minimal wrist drum-trigger: sensor and sample rate,
the feature you threshold for onset, how you map intensity to velocity, the channel and note numbers
for kick/snare/hi-hat, and the transport to the synth. Annotate each block with its latency budget so
the total lands under $15$ ms, and identify which block you expect to be the bottleneck.
\end{exercise}
